\documentclass[10pt,twocolumn]{article}

\usepackage[a4paper,margin=0.72in,columnsep=0.23in]{geometry}
\usepackage{times}
\usepackage{microtype}
\usepackage{graphicx}
\usepackage{booktabs}
\usepackage{array}
\usepackage{enumitem}
\usepackage{xcolor}
\usepackage{xurl}
\usepackage{listings}
\usepackage[hidelinks]{hyperref}

\definecolor{linkblue}{RGB}{0,76,151}
\definecolor{codebg}{RGB}{248,248,248}
\definecolor{codeframe}{RGB}{210,210,210}
\definecolor{slotbg}{RGB}{252,252,252}
\definecolor{slotline}{RGB}{180,180,180}

\hypersetup{
  colorlinks=true,
  linkcolor=linkblue,
  urlcolor=linkblue,
  citecolor=linkblue
}

\lstdefinestyle{bashstyle}{
  basicstyle=\ttfamily\scriptsize,
  backgroundcolor=\color{codebg},
  frame=single,
  rulecolor=\color{codeframe},
  breaklines=true,
  columns=fullflexible,
  keepspaces=true,
  showstringspaces=false,
  xleftmargin=0.4em,
  xrightmargin=0.4em,
  aboveskip=0.55em,
  belowskip=0.55em
}
\lstset{style=bashstyle}

\newcommand{\code}[1]{\nolinkurl{#1}}
\newcommand{\slotbox}[2]{%
  \noindent\fcolorbox{slotline}{slotbg}{%
    \begin{minipage}{0.94\linewidth}
      \vspace{0.35em}
      \textbf{#1.} #2
      \vspace{0.35em}
    \end{minipage}
  }%
}
\newcommand{\imageplaceholder}[1]{%
  \fcolorbox{slotline}{slotbg}{%
    \begin{minipage}[c][1.38in][c]{0.90\linewidth}
      \centering
      \textbf{Image slot}\\[0.35em]
      Replace this box with\\
      \footnotesize
      \texttt{\string\includegraphics[width=\string\linewidth]\{}%
      \nolinkurl{#1}%
      \texttt{\}}
    \end{minipage}
  }%
}

\title{\vspace{-0.6em}\bfseries TEDEX V1.3.0: A Reproducible Workspace for Thermoelectric Data Analysis}
\author{
  Heyang Chen\\
  Nanyang Techological University, Northwestern Thermoelectrics\\
  \texttt{te\_agent\_workspace}
}
\date{\vspace{-0.8em}}

\begin{document}
\maketitle

\begin{abstract}
TEDEX is a Python-based workspace for thermoelectric materials research. It
organizes raw ZEM/LFA/XRD measurements, sample metadata, processed transport
tables, publication-style figures, batch assessments, and local optimization
reports in a single reproducible project layout. This document gives a compact
system overview and a practical operating guide. It is intentionally formatted
like a short CS/AI systems paper: the first half describes the architecture and
dataflow, while the second half provides command-driven operation cards with
fillable slots for bash commands, terminal output, and generated figures.
\end{abstract}

\noindent\textbf{Keywords:}
thermoelectric materials; data pipeline; XRD; transport properties; reproducible
analysis; Bayesian optimization; scientific plotting

\section{Introduction}

Modern thermoelectric experiments produce heterogeneous data: instrument exports
from ZEM and LFA, XRD scans, sample ledgers, density and heat-capacity metadata,
plotting recipes, extracted literature tables, and downstream reports. TEDEX
provides a lightweight but structured workspace for turning those artifacts into
traceable analysis products.

The main design principle is simple: runnable entry points stay at the project
root, reusable implementation lives under \code{src/}, source data lives under
\code{data/}, and generated artifacts are separated into \code{results/} and
\code{outputs/}. This keeps interactive lab work fast while preserving enough
structure for batch reprocessing and paper-figure generation.

\section{System Overview}

Figure~\ref{fig:pipeline} summarizes the intended workflow. Raw experimental
files are synchronized into metadata ledgers, processed into sample-level CSV
tables, visualized by dedicated plotting tools, and optionally passed to
assessment or prediction modules.

\begin{figure}[t]
  \centering
  \fcolorbox{slotline}{slotbg}{%
    \begin{minipage}{0.92\linewidth}
      \centering
      \small
      \textbf{Raw data}\\
      \code{data/raw/CHY-xxxx/}\\[0.45em]
      $\downarrow$ metadata sync\\[0.35em]
      \textbf{Lab ledgers}\\
      \code{data/lab/*.json}, \code{lab_metadata.md}\\[0.45em]
      $\downarrow$ TE/XRD analysis\\[0.35em]
      \textbf{Processed tables and figures}\\
      \code{data/processed/}, \code{outputs/figures/}\\[0.45em]
      $\downarrow$ assessment and optimization\\[0.35em]
      \textbf{Reports}\\
      \code{results/reports/}, \code{results/assessments/}
    \end{minipage}
  }
  \caption{High-level TEDEX dataflow. Replace this schematic with a polished
  pipeline diagram if the document is used in a presentation or report.}
  \label{fig:pipeline}
\end{figure}

\begin{table}[t]
  \centering
  \caption{Top-level project layout.}
  \label{tab:layout}
  \begin{tabular}{@{}p{0.27\linewidth}p{0.62\linewidth}@{}}
    \toprule
    Path & Role \\
    \midrule
    \code{*.py} & Root-level runnable entry scripts. \\
    \code{src/} & Importable analysis, plotting, XRD, SPB, and agent modules. \\
    \code{scripts/} & Secondary CLI utilities for synchronization and export. \\
    \code{data/raw/} & Original instrument exports and unmodified source data. \\
    \code{data/processed/} & Cleaned transport-property CSV files. \\
    \code{data/lab/} & Batch ledgers, sample ledgers, and editable metadata. \\
    \code{configs/} & Plot recipes and reusable runtime configurations. \\
    \code{results/} & Reports, assessments, prediction tables, and payloads. \\
    \code{outputs/} & Figures, tables, logs, slide assets, and one-off exports. \\
    \bottomrule
  \end{tabular}
\end{table}

\section{Core Capabilities}

\paragraph{Transport analysis.}
The raw-data workflow reads paired ZEM/LFA files, calculates transport
properties, extracts compact sample features, and writes processed CSV outputs.
The main entry point is \code{python main.py analyze}.

\paragraph{TE plotting.}
\code{python main.py plot-te} renders single-property plots, combined summary figures, and
inter-batch comparisons. It supports axis limits, selected property panels,
multiple output formats, and no-preview batch rendering.

\paragraph{XRD plotting and lattice fitting.}
\code{python main.py plot-xrd} renders normalized XRD scans by default, supports
raw-count output with \code{--raw}, selected samples, optional PDF-card stick
patterns, and writes figures under
\code{outputs/figures/xrd/}. Lattice fitting utilities live under
\code{src/tools/xrd_lattice.py}.

\paragraph{Flexible plotting.}
\code{python main.py flexible} targets loosely structured CSV/TXT/XLSX files. It can
render line, scatter, bar, grouped-bar, multi-panel, and dual-axis figures from
direct CLI arguments or reusable JSON recipes.

\paragraph{Assessment and prediction.}
\code{python main.py assess} summarizes selected batches from extracted
feature payloads. \code{python main.py bayes} fits local Gaussian-process
models over lab or reference feature tables and ranks candidate compositions.

\paragraph{Agent-assisted interpretation.}
The \code{src/agents/} module adds model-assisted interpretation. It requires
the Python dependencies in \code{requirements.txt} and a private local API key
configured through \code{.env} or environment variables.

\section{Environment Setup}

All commands should be run from the workspace root:

\begin{lstlisting}[language=bash,caption={Workspace setup.}]
cd "/Users/chenheyang/Library/CloudStorage/OneDrive-NorthwesternUniversity/02-Northwestern Lab/te_agent_workspace"
pip install -r requirements.txt
\end{lstlisting}

Use the following commands to inspect the active interfaces:

\begin{lstlisting}[language=bash,caption={Entrypoint help commands.}]
python main.py analyze --help
python main.py plot-te --help
python main.py plot-xrd --help
python main.py flexible --help
python main.py assess --help
python main.py bayes --help
\end{lstlisting}

\section{Operating Cards}

Each operating card below is designed to be filled in after an actual run. The
\emph{Command} block records the exact bash command, the \emph{Terminal output}
block records the key log lines, and the \emph{Figure slot} shows where to place
the generated image.

\subsection{Metadata Synchronization}

Purpose: scan raw folders, update \code{data/lab/batches.json} and
\code{data/lab/samples.json}, and optionally round-trip the editable Markdown
ledger.

\begin{lstlisting}[language=bash,caption={Command slot: metadata sync.}]
# TODO: replace with the exact command used in the run.
python scripts/sync_lab_metadata.py --dry-run
python scripts/sync_lab_metadata.py
python scripts/sync_lab_markdown.py export
# Edit data/lab/lab_metadata.md, then import edits:
python scripts/sync_lab_markdown.py
\end{lstlisting}

\begin{lstlisting}[caption={Terminal-output slot: metadata sync.}]
# TODO: paste the key output lines here.
# Example:
# Found batch CHY-1048
# Updated data/lab/batches.json
# Updated data/lab/samples.json
\end{lstlisting}

\slotbox{Notes slot}{Record any manual metadata edits here, such as density,
\code{cp_value}, composition labels, annealing conditions, or known sample
issues.}

\subsection{Raw TE Analysis}

Purpose: process ZEM/LFA raw data for one or more batches and write cleaned
transport-property tables.

\begin{lstlisting}[language=bash,caption={Command slot: TE raw-data analysis.}]
# Single batch
python main.py analyze CHY-1048

# Multiple batches
python main.py analyze CHY-1038 CHY-1040

# Selected samples within one batch
python main.py analyze CHY-1054 --sample B C
\end{lstlisting}

\begin{lstlisting}[caption={Terminal-output slot: TE analysis.}]
# TODO: paste representative log lines here.
# Example:
# step1. Data processing
# Saved processed data to: data/processed/CHY-1048/...
# step2. SPB-derived transport columns
# Saved extracted features to: results/...
\end{lstlisting}

\begin{figure}[t]
  \centering
  \imageplaceholder{outputs/figures/te/CHY-1048/example.png}
  \caption{Figure slot for a representative TE output generated after raw-data
  analysis.}
  \label{fig:te-analysis-slot}
\end{figure}

\subsection{TE Plotting}

Purpose: generate single-property and combined thermoelectric figures.

\begin{lstlisting}[language=bash,caption={Command slot: TE plotting.}]
# Default single-property output
python main.py plot-te CHY-1040 --no-show

# Combined summary figure
python main.py plot-te CHY-1040 --plot-mode combined --no-show

# Inter-batch comparison with fixed axes
python main.py plot-te CHY-1038 CHY-1040 \
  --plot-mode combined \
  --comparison-id CHY1038_vs_CHY1040 \
  --xlim 300 800 \
  --no-show
\end{lstlisting}

\begin{lstlisting}[caption={Terminal-output slot: TE plotting.}]
# TODO: paste saved-file lines here.
# Example:
# Saved: outputs/figures/te/CHY-1040/seebeck.svg
# Saved: outputs/figures/te/inter-batch/CHY1038_vs_CHY1040/...
\end{lstlisting}

\begin{figure}[t]
  \centering
  \imageplaceholder{outputs/figures/te/CHY-1040/seebeck.png}
  \caption{Figure slot for a TE property plot, such as Seebeck, conductivity,
  power factor, or $zT$.}
  \label{fig:te-plot-slot}
\end{figure}

\subsection{XRD Plotting}

Purpose: render XRD scans for a batch, selected samples, or a direct XRD folder;
optionally overlay a cleaned PDF-card standard.

\begin{lstlisting}[language=bash,caption={Command slot: XRD plotting.}]
# Whole batch
python main.py plot-xrd CHY-1038 --no-show

# One normalized sample (default)
python main.py plot-xrd CHY-1038-A --no-show

# One raw-count sample
python main.py plot-xrd CHY-1038-A --raw --no-show

# Selected suffixes from an XRD folder
python main.py plot-xrd data/raw/CHY-1038/XRD \
  --select A C \
  --no-show

# With a PDF-card standard
python main.py plot-xrd CHY-1038 \
  --pdf-card CuInTe2_PDF_97_023_8958 \
  --no-show
\end{lstlisting}

\begin{lstlisting}[caption={Terminal-output slot: XRD plotting.}]
# TODO: paste saved-file lines or selected sample list here.
# Example:
# Loaded 3 XRD files
# Saved: outputs/figures/xrd/CHY-1038/...
\end{lstlisting}

\begin{figure}[t]
  \centering
  \imageplaceholder{outputs/figures/xrd/CHY-1038/example.png}
  \caption{Figure slot for a stacked or normalized XRD plot.}
  \label{fig:xrd-slot}
\end{figure}

\subsection{Flexible Plotting}

Purpose: render figures from ad hoc CSV/TXT/XLSX inputs or reusable JSON
recipes.

\begin{lstlisting}[language=bash,caption={Command slot: flexible plotting.}]
# Direct CLI mode
python main.py flexible \
  --data data/demo/flexible_plotting/literature_comparison_loose.csv \
  --kind scatter \
  --x "Ag fraction x" \
  --y "Best zT" \
  --xlabel "Ag fraction, $x$" \
  --ylabel "$ZT_{\\mathrm{max}}$" \
  --stem demo_scatter \
  --formats png pdf \
  --no-show

# Recipe mode
python main.py flexible \
  --recipe configs/flexible_plot_demos/temperature_seebeck_line.json \
  --no-show
\end{lstlisting}

\begin{lstlisting}[caption={Terminal-output slot: flexible plotting.}]
# TODO: paste generated files here.
# Example:
# outputs/figures/flexible_cli/demo_scatter.png
# outputs/figures/flexible_cli/demo_scatter_normalized.csv
\end{lstlisting}

\begin{figure}[t]
  \centering
  \imageplaceholder{outputs/figures/flexible_cli/demo_scatter.png}
  \caption{Figure slot for a flexible-plotting output.}
  \label{fig:flexible-slot}
\end{figure}

\subsection{Assessment and Bayesian Prediction}

Purpose: summarize extracted transport features and rank candidate compositions
for follow-up measurements.

\begin{lstlisting}[language=bash,caption={Command slot: assessment and prediction.}]
# Assess selected batches
python main.py assess CHY-1040 CHY-1048

# Assess every batch with extracted features
python main.py assess --all

# Local GP prediction on lab data
python main.py bayes --source lab --target ZT_max --top 10

# Reference-data prediction
python main.py bayes \
  --source reference \
  --reference-id GeSe \
  --target ZT_max \
  --pdf
\end{lstlisting}

\begin{lstlisting}[caption={Terminal-output slot: prediction.}]
# TODO: paste top candidates or report paths here.
# Example:
# Top candidate: ...
# Report: results/bayesian_predictions/...
# Figure: outputs/figures/bayesian_predictions/...
\end{lstlisting}

\section{Recommended End-to-End Protocol}

For a new experimental batch, use the following protocol:

\begin{enumerate}[leftmargin=*]
  \item Place raw files under \code{data/raw/<batch_id>/}. Keep original
  instrument exports unchanged.
  \item Run metadata synchronization and fill missing density, heat capacity,
  composition, and synthesis annotations.
  \item Process the raw TE data with \code{python main.py analyze}.
  \item Render TE figures and inspect axis limits.
  \item Render XRD figures and, when appropriate, compare with a PDF-card
  standard.
  \item Run batch assessment and record the main finding in \code{results/}.
  \item If enough structured data exists, run a Bayesian prediction pass and
  use the ranked candidates to plan follow-up samples.
\end{enumerate}

\begin{lstlisting}[language=bash,caption={End-to-end command template.}]
python scripts/sync_lab_metadata.py
python scripts/sync_lab_markdown.py export
# TODO: edit data/lab/lab_metadata.md
python scripts/sync_lab_markdown.py
python main.py analyze CHY-1048
python main.py plot-te CHY-1048 --seebeck --no-show
python main.py plot-xrd CHY-1048 --no-show
python main.py assess CHY-1048
\end{lstlisting}

\section{Reproducibility Notes}

Generated outputs should be stored in predictable locations. Standard TE plots
belong under \code{outputs/figures/te/}; XRD plots belong under
\code{outputs/figures/xrd/}; flexible one-off plots belong under
\code{outputs/figures/flexible_cli/} or \code{outputs/figures/flexible/};
assessment reports belong under \code{results/assessments/}; Bayesian reports
belong under \code{results/bayesian_predictions/}. Private credentials and
local-only configuration files should remain outside versioned or public
artifacts.

\section{Editable Example Slots}

Use the following minimal pattern when adding a new operation example. Copy the
three blocks together so each example records its command, textual output, and
visual output.

\begin{lstlisting}[language=bash,caption={Generic command slot.}]
# TODO: paste the exact bash command here.
python your_script.py --your-option value --no-show
\end{lstlisting}

\begin{lstlisting}[caption={Generic terminal-output slot.}]
# TODO: paste only the important output lines here.
# Keep this short enough to fit in one column.
\end{lstlisting}

\begin{figure}[t]
  \centering
  \imageplaceholder{path/to/generated_figure.png}
  \caption{Generic output-image slot. Replace the placeholder with an actual
  figure once the command has been run.}
  \label{fig:generic-slot}
\end{figure}

\section{Pointers}

More detailed command references are maintained in
\code{docs/COMMANDS.md}, \code{docs/FLEXIBLE_PLOTTING.md}, and
\code{docs/PROJECT_STRUCTURE.md}. When the documentation and script behavior
disagree, prefer the current \code{--help} output from the corresponding root
entry point.

\end{document}
